\documentclass[11pt,a4paper]{article}
\usepackage[margin=2cm]{geometry}
\usepackage{amsmath,amssymb}
\usepackage{xcolor}
\usepackage{enumitem}
\usepackage{tcolorbox}
\usepackage{array}
\tcbuselibrary{breakable,skins}

\setlength{\parindent}{0pt}
\setlength{\parskip}{4pt}

\definecolor{keybg}{HTML}{FFF3CD}
\definecolor{keyrule}{HTML}{B08800}
\definecolor{traybg}{HTML}{EEF4FA}
\definecolor{trayrule}{HTML}{4B6584}

\newcommand{\pred}[1]{\textsc{#1}}
\newcommand{\Result}{\pred{Result}}
\newcommand{\Poss}{\pred{Poss}}
\newcommand{\key}{\textbf{\textcolor{keyrule}{\large$\bigstar$}}\ }

\newtcolorbox{keybox}[1][]{
  enhanced, breakable,
  colback=keybg, colframe=keyrule,
  fonttitle=\bfseries, title={#1},
  left=6pt, right=6pt, top=4pt, bottom=4pt
}
\newtcolorbox{tray}[1][]{
  enhanced, breakable,
  colback=traybg, colframe=trayrule,
  fonttitle=\bfseries, title={#1},
  left=6pt, right=6pt, top=4pt, bottom=4pt
}

\title{\vspace{-1.5cm}AI Supo 2 --- Topics covered\\\large Q1, Q2 (Sec 6) and Q1 (Sec 7)}
\author{}
\date{}

\begin{document}
\maketitle
\vspace{-1cm}

\hfill \textbf{\textcolor{keyrule}{$\bigstar$}}\ = 80/20 essentials, learn cold.

\hrulefill

\section*{Section 6 --- Situation Calculus}

\subsection*{1.1\ \ The language}

Situation calculus encodes a changing world inside first-order logic.

\begin{keybox}[\key Core ontology]
\begin{tabular}{@{}p{3.3cm}p{12cm}@{}}
\textbf{Situation} $s$ & A snapshot of the world. $s_0$ is the initial situation. Situations are objects in the logic, so you can quantify over them.\\
\textbf{Fluent} & A predicate whose truth value depends on a situation. Last argument is always $s$.\ \ E.g.\ $\pred{At}(\text{Robot}, \text{kitchen}, s)$, $\pred{HaveArrow}(s)$.\\
\textbf{Action} & A first-class object (\emph{not} a situation). E.g.\ $\pred{Go}(\text{kitchen}, \text{hall})$, $\pred{Shoot}$.\\
\textbf{$\Result(a,s)$} & A function returning the situation that results from doing action $a$ in situation $s$.\\
\textbf{$\pred{Goal}(s)$} & A fluent meaning ``the goal is satisfied in $s$''.\\
\end{tabular}
\end{keybox}

Static facts about the world (e.g.\ $\pred{Adjacent}(p_1,p_2)$, $\pred{Wall}(p)$) are \emph{not} fluents --- they don't depend on $s$.

\subsection*{1.2\ \ Composing actions: \pred{Sequence}/\pred{Do}}

$\pred{Sequence}(a, s, s')$ means ``executing the action list $a$ from $s$ produces $s'$''. Defined recursively:
\[
  \pred{Sequence}([\,], s, s),
  \qquad
  \pred{Sequence}([a \mid \mathit{rest}], s, s') \iff \pred{Sequence}(\mathit{rest},\, \Result(a, s),\, s').
\]

\subsection*{1.3\ \ The frame problem}

When action $a$ happens, almost \emph{everything} stays the same --- only a tiny fraction changes. Stating ``everything else is unchanged'' naively requires $O(\text{actions} \times \text{fluents})$ frame axioms. That's the \textbf{frame problem}.

\begin{keybox}[\key Successor-state axioms (the fix)]
For \emph{each fluent} $F$, write \emph{one} axiom of the form
\[
  F(\vec{x}, \Result(a, s)) \iff \underbrace{\Phi^+(\vec{x}, a, s)}_{\text{$a$ makes $F$ true}}\ \vee\ \Big(F(\vec{x}, s)\ \wedge\ \neg \underbrace{\Phi^-(\vec{x}, a, s)}_{\text{$a$ makes $F$ false}}\Big).
\]
This combines effect axioms \emph{and} frame axioms in a single sentence per fluent. Reasoning becomes tractable.
\end{keybox}

\textbf{Common student error:} writing one axiom \emph{per action} rather than per fluent. Successor-state axioms close over \emph{all} actions $a$ at once.

\subsection*{1.4\ \ Possibility axioms}

For each action $A$, give one axiom $\Poss(A, s) \iff \ldots$ saying when $A$ is applicable in $s$. This is the precondition.

\begin{tray}[Recipe for every action]
\begin{enumerate}[leftmargin=*,itemsep=0pt]
  \item One \textbf{possibility} axiom: $\Poss(A,s) \iff \ldots$
  \item One \textbf{successor-state} axiom \emph{per fluent} that any action (including $A$) can affect.
\end{enumerate}
\end{tray}

\subsection*{1.5\ \ Planning queries (Q1)}

There are \emph{two} ways to phrase ``find me a plan'' as a logical query.

\begin{keybox}[\key Form A --- explicit situation]
\[
  \exists a\, \exists s\, .\ \pred{Sequence}(a, s_0, s) \,\wedge\, \pred{Goal}(s)
\]
Read: ``find an action list $a$ and a situation $s$ such that running $a$ from $s_0$ produces $s$, and the goal holds in $s$.'' A constructive proof binds $a$ to a plan.
\end{keybox}

\begin{keybox}[\key Form B --- situation hidden inside the action list]
\[
  \exists \mathit{actionList}\, .\ \pred{Goal}(\dots \mathit{actionList} \dots)
\]
No explicit $s$. The goal predicate must now take an action list as argument.
\end{keybox}

\textbf{How to make Form B work (the Q1 answer in three moves):}

\begin{enumerate}[leftmargin=*,itemsep=1pt]
  \item Define an auxiliary $\pred{Do}$ that folds the action list through $\Result$ starting from $s_0$:
  \[
    \pred{Do}([\,], s) = s, \qquad \pred{Do}([a \mid \mathit{rest}], s) = \pred{Do}(\mathit{rest},\, \Result(a, s)).
  \]
  \item Redefine the goal predicate so the situation is \emph{internal} to it:
  \[
    \pred{Goal}(\mathit{actionList}) \iff \pred{GoalHolds}(\pred{Do}(\mathit{actionList}, s_0)).
  \]
  Here $\pred{GoalHolds}(s)$ is the original situation-based goal fluent.
  \item Equivalence with Form A: the existentially quantified $s$ is exactly $\pred{Do}(\mathit{actionList}, s_0)$. So Form B holds iff Form A would have with $s$ instantiated to that value. The same successor-state axioms make either form evaluable.
\end{enumerate}

\textbf{Why bother?} The witness from a constructive proof of Form B is \emph{just} the action list --- which is what the planner actually wants. Situations are a derived intermediate quantity that Form A forces you to quantify over.

\subsection*{1.6\ \ Wumpus World: Shoot, Move, Orient (Q2)}

\textbf{Constants \& helpers (assumed non-fluent unless marked):}
\begin{itemize}[leftmargin=*,itemsep=0pt]
  \item Locations $[x,y]$; orientations $\in \{0, 90, 180, 270\}$; \pred{Agent}, \pred{Wumpus} constants.
  \item $\pred{LocationToward}([x,y], d)$: the adjacent cell in direction $d$.
  \item $\pred{LocationAhead}(s) = \pred{LocationToward}(\pred{loc}(s),\pred{ori}(s))$.
  \item $\pred{Wall}(p)$, $\pred{Adjacent}(p_1, p_2)$: static map predicates.
\end{itemize}

\textbf{Fluents:} $\pred{At}(\pred{Agent},[x,y],s)$, $\pred{Orientation}(\pred{Agent},d,s)$, $\pred{HaveArrow}(s)$, $\pred{WumpusAlive}(s)$.

\textbf{Helper fluent (line-of-fire):}
\[
  \pred{FacingWumpus}(s) \iff \exists [x,y].\ \pred{At}(\pred{Wumpus},[x,y],s)\,\wedge\,\pred{InLineOfFire}(\pred{Agent},[x,y],s).
\]

\begin{keybox}[\key Shoot]
\textbf{Possibility:} $\Poss(\pred{Shoot}, s) \iff \pred{HaveArrow}(s)$.\\[2pt]
\textbf{Successor-state for \pred{HaveArrow}:}
$\pred{HaveArrow}(\Result(a,s)) \iff \pred{HaveArrow}(s)\,\wedge\,a \neq \pred{Shoot}$.\\[2pt]
\textbf{Successor-state for \pred{WumpusAlive}:}
\[
  \pred{WumpusAlive}(\Result(a,s)) \iff \pred{WumpusAlive}(s)\,\wedge\,\neg\big(a=\pred{Shoot}\,\wedge\,\pred{FacingWumpus}(s)\big).
\]
\end{keybox}

\begin{keybox}[\key Move (Forward)]
\textbf{Possibility:} $\Poss(\pred{Forward}, s) \iff \neg \pred{Wall}(\pred{LocationAhead}(s))$.\\[2pt]
\textbf{Successor-state for \pred{At}:}
\begin{align*}
  \pred{At}(\pred{Agent}, p, \Result(a, s)) \iff{}&
  \big(a = \pred{Forward}\,\wedge\,p = \pred{LocationAhead}(s)\big)\\
  \vee{}& \big(a \neq \pred{Forward}\,\wedge\,\pred{At}(\pred{Agent}, p, s)\big).
\end{align*}
\end{keybox}

\begin{keybox}[\key Orient (TurnLeft, TurnRight)]
$\pred{LeftTurn}(d) = (d+90) \bmod 360,\quad \pred{RightTurn}(d) = (d-90) \bmod 360$.\\[2pt]
\textbf{Possibility:} $\Poss(\pred{TurnLeft},s) \iff \top$, $\ \Poss(\pred{TurnRight},s) \iff \top$.\\[2pt]
\textbf{Successor-state for \pred{Orientation}:}
\begin{align*}
\pred{Orientation}(\pred{Agent},d,\Result(a,s)) \iff{}
& (a = \pred{TurnLeft}\,\wedge\,d = \pred{LeftTurn}(\pred{ori}(s)))\\
\vee{}& (a = \pred{TurnRight}\,\wedge\,d = \pred{RightTurn}(\pred{ori}(s)))\\
\vee{}& (a \neq \pred{TurnLeft}\,\wedge\,a \neq \pred{TurnRight}\\
& \quad \wedge\,\pred{Orientation}(\pred{Agent},d,s)).
\end{align*}
\end{keybox}

\subsection*{1.7\ \ Common pitfalls (Sec 6)}

\begin{itemize}[leftmargin=*,itemsep=0pt]
  \item Writing effect axioms only --- forgetting the frame part. Use successor-state form so frames are automatic.
  \item Writing one successor-state axiom per action instead of per fluent.
  \item Treating $\pred{Wall}$, $\pred{Adjacent}$ as fluents and writing pointless successor-state axioms for them.
  \item Forgetting the outer universal $\forall a, s$ when writing the axioms formally.
\end{itemize}

\hrulefill

\section*{Section 7 --- Heuristics for Planning}

\subsection*{2.1\ \ Planning vs search}

\begin{keybox}[\key The four differences (worth knowing cold)]
\begin{tabular}{@{}p{3.4cm}p{6cm}p{6cm}@{}}
& \textbf{Search} & \textbf{Planning}\\
\hline
State representation & Opaque (black box) & Structured (set of fluents)\\
Actions & Arbitrary transitions & Precondition + add-list + delete-list\\
Goal test & Single function & Logical condition over fluents\\
Heuristic source & Domain-specific, supplied by human & Can be \textbf{derived from the structure}\\
\end{tabular}
\end{keybox}

\subsection*{2.2\ \ Why planning gives heuristics ``for free''}

In ordinary search the human writes the heuristic. In planning, each operator's preconditions and add/delete lists are visible to the planner --- so the planner can \emph{automatically} construct a relaxed problem and use its optimal cost as $h$.

\begin{keybox}[\key The recipe]
\textbf{Drop part of the problem to get an easier (``relaxed'') problem whose optimal cost is a lower bound on the original.} Solving the relaxation must be substantially cheaper (preferably polynomial) and the bound must be admissible.
\end{keybox}

\subsection*{2.3\ \ Relaxation heuristics}

\textbf{(a) Ignore delete lists (delete-relaxation).} Drop every action's delete-list. Nothing ever becomes false; literals only accumulate. The relaxed problem is monotone and solvable in polynomial time. Its length is $h^+$, the canonical admissible heuristic. \textbf{Used by FastForward (FF) --- one of the most successful planners.}

\textbf{(b) Ignore preconditions.} Drop all preconditions. Every action always applies. Optimal relaxed-plan length = minimum number of actions whose add-lists cover the goal (a set-cover problem). Looser than $h^+$ but easy to compute.

\textbf{(c) Subgoal independence.} Plan for each goal conjunct in isolation, then combine:
\begin{itemize}[leftmargin=*,itemsep=0pt]
  \item $h_{\max}(s) = \max_g \mathrm{cost}(g)$ \ --- admissible (under-counts).
  \item $h_{\text{add}}(s) = \sum_g \mathrm{cost}(g)$ \ --- not admissible (double-counts) but often more informative.
\end{itemize}

\subsection*{2.4\ \ Other useful heuristics}

\textbf{Number of unsatisfied goal conditions remaining.} Quick and dirty; not admissible in general but cheap and surprisingly effective on simple problems.

\textbf{Planning-graph level.} The depth at which a goal first becomes reachable in a planning graph is an admissible heuristic (deeper graph = harder problem). Used by Graphplan.

\textbf{Helpful actions.} Restrict the branching factor at each search node to only those actions that appear in the relaxed plan from that node. Trades completeness for speed.

\subsection*{2.5\ \ How to structure the Q1 (Sec 7) answer}

A strong answer names at least \textbf{three heuristics from different families} and links them back to the question. A good template:

\begin{enumerate}[leftmargin=*,itemsep=1pt]
  \item State the key idea: planning's structured representation lets heuristics be derived \emph{automatically} from the problem, unlike generic search.
  \item Give the relaxation recipe (drop part of the problem; lower bound on cost).
  \item Name two or three concrete heuristics: at least one relaxation-based (ignore deletes / ignore preconditions / subgoal independence), and at least one structural (planning-graph level / unsatisfied-goals count / helpful actions).
  \item Close by linking back: planning's edge over generic search is precisely this ability to compute its own heuristic.
\end{enumerate}

\hrulefill

\section*{The 80/20 cheat-sheet (memorise this page)}

\begin{tray}[Sec 6 --- non-negotiable]
\begin{itemize}[leftmargin=*,itemsep=0pt]
  \item Situations, fluents, actions, $\Result(a,s)$ --- definitions on the tip of your tongue.
  \item For every action: \textbf{one possibility axiom, one successor-state axiom per affected fluent.}
  \item Successor-state schema: $F(\vec{x}, \Result(a, s)) \iff \Phi^+\,\vee\,(F(\vec{x},s) \wedge \neg \Phi^-)$.
  \item Two query forms: $\exists a\,\exists s.\ \pred{Sequence}(a,s_0,s) \wedge \pred{Goal}(s)$ vs $\exists \mathit{actionList}.\ \pred{Goal}(\mathit{actionList})$.
  \item Form B works by defining $\pred{Goal}(\mathit{actionList}) \iff \pred{GoalHolds}(\pred{Do}(\mathit{actionList}, s_0))$.
\end{itemize}
\end{tray}

\begin{tray}[Sec 7 --- non-negotiable]
\begin{itemize}[leftmargin=*,itemsep=0pt]
  \item Planning vs search: structured state, structured actions, derivable heuristics.
  \item The relaxation recipe: drop part of the problem; cost of relaxed solution = lower bound on real cost.
  \item Two relaxations you must know: \textbf{ignore deletes} ($h^+$, FastForward), \textbf{ignore preconditions} (set cover).
  \item Subgoal-independence: $h_{\max}$ admissible, $h_{\text{add}}$ informative.
\end{itemize}
\end{tray}

\end{document}
